\chapter{Minimal Code Recipes}
\label{app:minimal-code-recipes}

\textit{These recipes are pseudocode-level implementation plans. They are not project files and they intentionally omit exact source-paper parameters until the coordinator verifies citations and parameter tables. Their purpose is to make the research workflow unambiguous before writing production simulation code.}

% ============================================================
\section{Recipe 1: Build E-Cluster Connectivity}
\label{app:recipe-build-e-cluster-connectivity}
% ============================================================

The first recipe creates a clustered excitatory connectivity matrix while keeping mean recurrent excitation controlled.

\begin{lstlisting}[language=Python,caption={Pseudocode for E-cluster connectivity.}]
inputs:
    Q                    # number of excitatory clusters
    n_E                  # excitatory neurons per cluster
    N_I                  # inhibitory neurons
    p_EE_base            # baseline E->E connection probability
    g_EE                 # within-cluster strength multiplier
    p_EI, p_IE, p_II     # remaining connection probabilities
    seed

set random seed
assign each excitatory neuron i to cluster c[i] in {1,...,Q}

p_EE_plus  = g_EE * p_EE_base
p_EE_minus = (Q * p_EE_base - p_EE_plus) / (Q - 1)

for each ordered E target i and E source j:
    if i == j:
        W_EE[i,j] = 0
    else if c[i] == c[j]:
        W_EE[i,j] = Bernoulli(p_EE_plus) * J_EE
    else:
        W_EE[i,j] = Bernoulli(p_EE_minus) * J_EE

draw E<-I, I<-E, and I<-I connections from their baseline rules
store cluster labels, probabilities, weights, and seed
\end{lstlisting}

The normalization line is only valid when it keeps $p_{EE}^{-}$ in $[0,1]$. If it does not, switch from probability clustering to weight clustering or choose a smaller $g_{EE}$.

% ============================================================
\section{Recipe 2: Simulate the Minimal LIF Baseline}
\label{app:recipe-simulate-lif-baseline}
% ============================================================

The microscopic simulation should produce spike times and population summaries in one pass.

\begin{lstlisting}[language=Python,caption={Pseudocode for a current-based LIF simulation.}]
inputs:
    connectivity matrices
    neuron parameters: tau_m, V_rest, V_th, V_reset, tau_ref
    synapse parameters: tau_s, J values
    external drive I_ext
    time step dt and duration T

initialize V_i near V_rest
initialize synaptic traces y_i = 0
initialize refractory counters

for each time step t:
    decay all synaptic traces:
        y_i <- y_i + dt * (-y_i / tau_s)

    compute synaptic current for every neuron:
        I_i = I_ext_i + sum_E W_E[i,j] y_j - sum_I W_I[i,j] y_j

    update non-refractory membrane potentials:
        V_i <- V_i + dt/tau_m * (-(V_i - V_rest) + R_m * I_i)

    for neurons crossing threshold:
        record spike time
        set V_i = V_reset
        set refractory counter = tau_ref
        add spike impulse to y_i

    decrement refractory counters

save:
    spike times
    cluster labels
    binned E and I rates
    optional E/I current summaries
\end{lstlisting}

The simulation should be tested first with clustering disabled. If the unclustered network is not active, balanced, and irregular, the clustered version will be difficult to interpret.

% ============================================================
\section{Recipe 3: Detect Cluster States}
\label{app:recipe-detect-cluster-states}
% ============================================================

Cluster-state detection converts spikes into the variables used in the mesoscopic theory.

\begin{lstlisting}[language=Python,caption={Pseudocode for cluster-state detection.}]
inputs:
    spike times
    cluster labels
    smoothing kernel kappa
    theta_on, theta_off
    minimum_duration

for each cluster k:
    r_hat[k,t] = smoothed spike count of neurons in cluster k / n_E

for each cluster k:
    state[k,0] = 0
    for each time bin t:
        if r_hat[k,t] > theta_on:
            state[k,t] = 1
        else if r_hat[k,t] < theta_off:
            state[k,t] = 0
        else:
            state[k,t] = state[k,t-1]

remove active intervals shorter than minimum_duration
compute:
    active_cluster_count[t] = sum_k state[k,t]
    lifetimes for each active interval
    transition times
\end{lstlisting}

Thresholds should be chosen once per analysis protocol, preferably from the low/high valley of the pooled cluster-rate distribution. Do not tune thresholds separately for every parameter point.

% ============================================================
\section{Recipe 4: Fit or Choose Mesoscopic Transfer Functions}
\label{app:recipe-fit-transfer-functions}
% ============================================================

The transfer function maps input statistics to expected population rate. A first project can choose a smooth nonlinearity; a stronger project estimates it from microscopic simulations.

\begin{lstlisting}[language=Python,caption={Pseudocode for transfer-function estimation.}]
inputs:
    isolated or embedded population model
    grid of mean input values mu_grid
    optional grid of input variances sigma_grid
    simulation duration per grid point

for each mu in mu_grid:
    for each sigma in sigma_grid:
        drive population with fluctuating input having mean mu and variance sigma
        simulate until transient has passed
        estimate steady firing rate r_bar
        store data point (mu, sigma, r_bar)

fit Phi(mu, sigma) with a smooth constrained curve:
    nonnegative output
    monotone in mu when appropriate
    bounded or saturating if needed

validate:
    predicted rates match held-out input points
    gain Phi' is not unrealistically steep
\end{lstlisting}

If the final mesoscopic model only uses $\Phi(u)$, the variance dimension can be collapsed. If refractory-density theory is used, the transfer function may be replaced by age-structured dynamics instead.

% ============================================================
\section{Recipe 5: Integrate the E-Only Mesoscopic Model}
\label{app:recipe-integrate-e-only-mesoscopic}
% ============================================================

This recipe tests whether the reduced E-only model can reproduce cluster switching.

\begin{lstlisting}[language=Python,caption={Pseudocode for E-only mesoscopic integration.}]
inputs:
    Q, n_E
    tau_E, tau_I
    J_plus, J_minus, J_EI, J_IE
    transfer functions Phi_E, Phi_I
    finite-size noise function sigma_E
    dt, T, seed

initialize r_E[k] near low-rate state
initialize r_I

for each time step:
    u_E[k] = h_E[k] + J_plus*r_E[k]
             + J_minus*sum_{l != k} q[l]*r_E[l]
             - J_EI*r_I
    u_I = h_I + J_IE*sum_l q[l]*r_E[l]

    drift_E[k] = (-r_E[k] + Phi_E(u_E[k])) / tau_E
    drift_I    = (-r_I    + Phi_I(u_I))    / tau_I

    r_E[k] <- r_E[k] + dt*drift_E[k]
              + sqrt(dt) * sigma_E(r_E[k]) / sqrt(n_E) * normal()
    r_I    <- r_I + dt*drift_I

    enforce nonnegative rates by reducing dt or applying documented truncation

save r_E, r_I, active-cluster indicators, and summary statistics
\end{lstlisting}

Always run the same parameter point with the noise amplitude set to zero. That determines whether switching is finite-size driven or already present in the deterministic equations.

% ============================================================
\section{Recipe 6: Integrate the E/I Mesoscopic Extension}
\label{app:recipe-integrate-ei-mesoscopic}
% ============================================================

The E/I extension uses two variables per cluster and four coupling matrices.

\begin{lstlisting}[language=Python,caption={Pseudocode for E/I mesoscopic integration.}]
inputs:
    Q, n_E, n_I
    W_EE, W_EI, W_IE, W_II
    h_E[k], h_I[k]
    tau_E, tau_I
    Phi_E, Phi_I
    sigma_E, sigma_I
    dt, T, seed

initialize r_E[k] and r_I[k]

for each time step:
    for each cluster k:
        u_E[k] = h_E[k]
                 + sum_l W_EE[k,l] * r_E[l]
                 - sum_l W_EI[k,l] * r_I[l]

        u_I[k] = h_I[k]
                 + sum_l W_IE[k,l] * r_E[l]
                 - sum_l W_II[k,l] * r_I[l]

        drift_E[k] = (-r_E[k] + Phi_E(u_E[k])) / tau_E
        drift_I[k] = (-r_I[k] + Phi_I(u_I[k])) / tau_I

        r_E[k] <- r_E[k] + dt*drift_E[k]
                  + sqrt(dt)*sigma_E(r_E[k])/sqrt(n_E)*normal()
        r_I[k] <- r_I[k] + dt*drift_I[k]
                  + sqrt(dt)*sigma_I(r_I[k])/sqrt(n_I)*normal()

save:
    r_E, r_I
    local balance residuals
    E/I tracking correlations
    state occupancies and lifetimes
\end{lstlisting}

The first comparison should change only the E/I-structure parameters relative to the E-only baseline. Retuning many old parameters at the same time makes the mechanism hard to identify.

% ============================================================
\section{Recipe 7: Fixed Points and Stability}
\label{app:recipe-fixed-points-stability}
% ============================================================

Fixed-point analysis turns a simulation regime into a dynamical-systems claim.

\begin{lstlisting}[language=Python,caption={Pseudocode for fixed-point and stability analysis.}]
inputs:
    deterministic drift function f(x; theta)
    candidate initial states
    numerical solver tolerance

fixed_points = empty list

for each initial state x0:
    solve f(x; theta) = 0 starting from x0
    if solution is valid and not already in fixed_points:
        add solution

for each fixed point x_star:
    compute Jacobian J = derivative of f at x_star
    compute eigenvalues lambda
    classify:
        stable if max(real(lambda)) < 0
        saddle if some real parts positive and some negative
        near-critical if max(real(lambda)) is close to 0

return fixed point table:
    state pattern
    rates
    leading eigenvalue
    dominant eigenvector interpretation
\end{lstlisting}

For clustered systems, initialize the solver near no-active, one-active, and multiple-active patterns. Otherwise the solver may miss fixed points with small basins.

% ============================================================
\section{Recipe 8: Cluster-Size Scaling Test}
\label{app:recipe-cluster-size-scaling}
% ============================================================

This is the decisive test for finite-size-driven switching.

\begin{lstlisting}[language=Python,caption={Pseudocode for cluster-size scaling.}]
inputs:
    base deterministic parameters theta
    cluster_sizes = [n1, n2, n3, ...]
    seeds per size

for each n in cluster_sizes:
    set noise amplitude proportional to 1/sqrt(n)
    keep deterministic drift f(x; theta) unchanged

    for each seed:
        simulate stochastic mesoscopic model
        detect cluster states with fixed detection rule
        compute dwell times and within-state variance

    summarize over seeds:
        mean lifetime
        median lifetime
        confidence interval or bootstrap interval
        within-state variance

plot:
    log(mean lifetime) versus n
    within-state variance versus 1/n
\end{lstlisting}

If microscopic simulations are used, preserving the deterministic drift while changing $n$ is harder because connectivity normalization must be adjusted carefully. Record the normalization protocol.

% ============================================================
\section{Recipe 9: Phase-Diagram Sweep}
\label{app:recipe-phase-diagram-sweep}
% ============================================================

The phase diagram combines deterministic and stochastic classifiers.

\begin{lstlisting}[language=Python,caption={Pseudocode for a phase-diagram sweep.}]
inputs:
    axis_1_values
    axis_2_values
    base parameters
    seeds
    deterministic classifier
    stochastic classifier

for value_1 in axis_1_values:
    for value_2 in axis_2_values:
        theta = update(base parameters, value_1, value_2)

        deterministic_result = analyze_fixed_points(theta)
        optional_lyapunov = estimate_largest_lyapunov(theta)

        stochastic_summaries = []
        for seed in seeds:
            trajectory = simulate_with_noise(theta, seed)
            summaries = compute_observables(trajectory)
            stochastic_summaries.append(summaries)

        label = classify_regime(deterministic_result,
                                optional_lyapunov,
                                stochastic_summaries)

        store label and scalar overlays

plot categorical labels
overlay leading eigenvalue, log lifetime, or Fano factor
attach representative trajectories for selected grid points
\end{lstlisting}

The classifier should mark ambiguous points explicitly. Ambiguity near bifurcation boundaries is information, not a plotting failure.

% ============================================================
\section{Recipe 10: Microscopic--Mesoscopic Comparison}
\label{app:recipe-micro-meso-comparison}
% ============================================================

The final recipe aligns the microscopic and mesoscopic outputs.

\begin{lstlisting}[language=Python,caption={Pseudocode for microscopic--mesoscopic comparison.}]
inputs:
    microscopic spike data and cluster labels
    mesoscopic rate trajectories
    common smoothing kernel
    common state detection rule

from microscopic data:
    compute r_hat_E[k,t] and r_hat_I[k,t]
    detect states
    compute occupancy, lifetimes, Fano factors, autocorrelations

from mesoscopic data:
    use r_E[k,t] and r_I[k,t]
    detect states with the same rule
    compute occupancy, lifetimes, rate autocorrelations
    optionally sample point-process counts from rates for Fano comparison

build comparison table:
    observable
    microscopic value
    mesoscopic value
    discrepancy
    likely explanation

rank discrepancies:
    transfer-function mismatch
    coupling mismatch
    noise-scale mismatch
    missing slow variable
    detection artifact
\end{lstlisting}

The comparison should end with a model revision decision. A mismatch is useful only if it points to the next modification or to a limit of the mesoscopic reduction.
